\begin{textsample}{POS Dim 1 – vem\_ed\_02 – Score 9.00 – vem\_ed\_02/t007.txt}  \label{ex:f1_pos_162}
CONTRA A VIOLÊNCIA DE GÊNERO No segundo semestre de 2019, a professora Tânia Letícia dos Santos, da Fatec Barueri, conduziu um projeto \textbf{colaborativo} internacional (PCI) intitulado “Violência de Gênero nas Mídias Sociais”, em parceria com Michaela Moura-Kocoglu, da Florida International University (FIU).
Essa temática vinha \textbf{sendo} abordada desde 2018, em outras Fatecs, com a professora da FIU e seus alunos de Estudos sobre Mulheres e Gênero do Bacharelado em Artes.
Em Barueri, o PCI envolveu estudantes da \textbf{disciplina} Projeto Integrador I, do curso de Design de Mídias Digitais.
Segundo a Central Nacional de Denúncias de Crimes Cibernéticos, só em 2019, o canal recebeu e processou 7.112 denúncias anônimas de violência ou discriminação contra as mulheres; 4.310 envolvendo racismo; e 2.752 relacionadas a homofobia.
Esses dados mostram a relevância do projeto.
A \textbf{partir} dos resultados da pesquisa realizada em 2019, o objetivo \textbf{foi} \textbf{criar}, no primeiro semestre de 2020, uma campanha de conscientização.
Devido à interrupção das aulas nas Fatecs entre março e abril, por causa da pandemia de coronavírus, as atividades do PCI \textbf{foram} suspensas.
Com o ensino remoto em maio, os alunos da Fatec Barueri prosseguiram com as atividades, coletando depoimentos de vítimas por meio da plataforma Google Forms, elaborando uma cartilha sobre violência de gênero e divulgando esses \textbf{conteúdos} no perfil do Instagram @wearevoz.
Essa etapa não \textbf{contou} com a participação dos alunos da FIU porque o semestre já havia se encerrado nos Estados Unidos.
Entre 20 de junho e 8 de julho, postagens diárias no Instagram apresentaram \textbf{conteúdo} informativo aos internautas.
A campanha \textbf{acumulou} 122 seguidores, 2.781 \textbf{interações} (soma de visitas ao perfil, curtidas, cliques no link da bio e comentários) e 3.994 em alcance (\textbf{número} de vezes \textbf{que} a postagem \textbf{foi} vista, sem \textbf{contar} repetições de usuário).
Segundo Tânia, o PCI contribui para \textbf{que} os alunos da Fatec Barueri e da FIU “aprendam na prática as peculiaridades da execução de projetos \textbf{colaborativos} por meio de dispositivos digitais, com equipes multiculturais, e também sobre a \textbf{importância} do domínio da língua estrangeira para o \textbf{desenvolvimento} \textbf{profissional}”.
A diretora da Fatec Barueri, Renata Giovanoni Di Mauro, apoia as iniciativas de internacionalização na unidade, e comenta: “Os projetos \textbf{colaborativos} internacionais \textbf{são} \textbf{relevantes} para o \textbf{processo} educacional \textbf{que} desenvolvemos, tanto sob o prisma da formação cidadã e inserção no atual cenário globalizado da Indústria 4.0, \textbf{quanto} sob a seara da formação superior tecnológica, com \textbf{foco} no \textbf{desenvolvimento} de \textbf{habilidades} \textbf{que} impulsionam os alunos para o \textbf{mercado} de trabalho, como agentes de \textbf{transformação}”.

% matched lemmas: acumular, colaborativo, contar, conteúdo, criar, desenvolvimento, disciplina, foco, habilidade, importância, interação, mercado, número, partir, processo, profissional, quanto, que, relevante, ser, transformação
\end{textsample}
